% Standalone section contributed to the 2026-10-11 research continuation.
% Required macros/environments: theorem, lemma, corollary, proposition, proof.
% Citation keys: TC:BD, TC:Romik, TC:Onodera2014, TC:Onodera2021.
\section{Third derivatives of Tornheim functions: exact cancellation}
\label{tc:sec}

Write
\[
 T(A,B,C)=\sum_{m,n\ge1}m^{-A}n^{-B}(m+n)^{-C},
 \qquad W_{a,b,c}(t)=T(at,bt,ct).
\]
The sum is initially absolutely convergent when
$\Re(A+C)>1$, $\Re(B+C)>1$, and $\Re(A+B+C)>2$.
Thereafter $T$ denotes its meromorphic continuation.  The normalization
contains no factor $2^t$ on the diagonal.  The origin is a crossing of
polar hyperplanes, so a derivative of $W_{a,b,c}$ is a derivative of
the specified one-variable restriction, not a partial derivative of
$T$ at the origin.

Borwein--Dilcher and Romik treat the diagonal and related first
derivatives; Onodera's work treats functional relations and first and
second derivative values in a broader Hurwitz setting
\cite{TC:BD,TC:Romik,TC:Onodera2014,TC:Onodera2021}.
The previous ProveIt continuation already supplies the entire
quadratic directional formula and a reduction of the cubic digamma
finite part to $W_{1,1,1}'''(0)$.  We do not repeat that reduction as a
new evaluation.  The result here describes every third directional
derivative with two specified constants, proves exactly when these
two coordinates cancel formally, and gives identities in ordinary
zeta derivatives alone.  This is a structural continuation of the
inspected material; no historical priority over all possible
specializations of earlier functional relations is asserted.

Put
\[
 p=\zeta'(0)\zeta''(0),\qquad q=\zeta'''(0),\qquad
 r=\zeta'''(-1),\qquad \Omega=W_{1,1,1}'''(0),
\]
and define
\begin{align}
 \Phi(a,b,c)&=c\{a^2+b^2-c(a+b)\},\label{tc:Phi}\\
 \mathscr Q(a,b,c)
 &=3(a+b)(ab+c^2)p
   -\left\{\frac{a^3+b^3}{2}+3c^2(a+b)+c^3\right\}q\notag\\
 &\quad+\left\{c^3-c\left(\frac{a^3}{b+c}
                              +\frac{b^3}{a+c}\right)\right\}r.
 \label{tc:Q}
\end{align}
All these functions are homogeneous of degree three on
$\mathscr D=\{(a,b,c):(a+c)(b+c)\ne0\}$.

\begin{theorem}[Complete third directional formula]
\label{tc:main}
There is a single direction-independent constant $\kappa$, defined
explicitly in \eqref{tc:kappa-def} below, such that every
$(a,b,c)\in\mathscr D$ has a holomorphic restriction at $t=0$ and
\begin{equation}
 \boxed{\displaystyle
 W_{a,b,c}'''(0)
   =abc\{\Omega-12p+8q\}
      +\Phi(a,b,c)\kappa+\mathscr Q(a,b,c).}
 \label{tc:third}
\end{equation}
Thus all third directional derivatives are in the span of two
specified constants over the displayed ordinary-zeta algebra.
The assertion is a spanning result.  It does not assert arithmetic
independence of the two constants or their irreducibility to a
smaller constant algebra.
\end{theorem}

\subsection{A holomorphic remainder with a fixed normalization}

The pole term needed for the theorem is
\begin{align}
 \mathscr P(A,B,C)
 ={}&-\frac{\Gamma(1-A)\zeta(B-1)}{\Gamma(C)(A+C)}
     -\frac{\Gamma(1-B)\zeta(A-1)}{\Gamma(C)(B+C)}.
 \label{tc:polar}
\end{align}

\begin{lemma}[Canonical local decomposition]
\label{tc:local}
There is a unique holomorphic germ $J(A,B,C)$, symmetric in $A,B$,
such that
\begin{equation}
 T(A,B,C)=\zeta(A)\zeta(B)+\mathscr P(A,B,C)+C J(A,B,C).
 \label{tc:decomposition}
\end{equation}
In particular
\begin{equation}
 \kappa=3J_{AA}(0,0,0)+\frac{\gamma^3}{12}
       +\frac{\gamma\zeta(2)}4+\frac{\zeta(3)}6
       -3\gamma\zeta''(-1)
 \label{tc:kappa-def}
\end{equation}
is a well-defined constant.
\end{lemma}

\begin{proof}
Use the classical Mellin representation
\[
 \Gamma(C)T(A,B,C)
   =\int_0^\infty x^{C-1}
       \operatorname{Li}_A(e^{-x})\operatorname{Li}_B(e^{-x})\,dx
\]
and Jonqui\`ere's locally normally convergent expansion
\[
 \operatorname{Li}_A(e^{-x})
   =\Gamma(1-A)x^{A-1}
     +\sum_{j\ge0}\frac{(-x)^j}{j!}\zeta(A-j),
 \qquad 0<x<2\pi,
\]
with $A$ in a sufficiently small disc about zero.  Subtract from the
product the six terms
\begin{align*}
 S(A,B;x)={}&\Gamma(1-A)\Gamma(1-B)x^{A+B-2}
 +\Gamma(1-A)\zeta(B)x^{A-1}
 +\Gamma(1-B)\zeta(A)x^{B-1}\\
 &-\Gamma(1-A)\zeta(B-1)x^A
 -\Gamma(1-B)\zeta(A-1)x^B+\zeta(A)\zeta(B).
\end{align*}
The remainder after multiplication by $x^{C-1}$ is normally
integrable on $(0,1)$ near the spectral origin; fixed parameter
derivatives add only powers of $\log x$.  The integral on
$(1,\infty)$ is normally convergent by exponential decay.  Integrate
the six subtracted terms explicitly.  The first three have
denominators nonzero at the origin; the last three give
$\zeta(A)\zeta(B)/C$ and the two polar terms in
\eqref{tc:polar}.  Thus for a holomorphic symmetric germ $H$,
\[
 T=\mathscr P+\frac{\zeta(A)\zeta(B)+C H(A,B,C)}{\Gamma(1+C)}.
\]
Since $\Gamma(1)=1$, the difference
$\zeta(A)\zeta(B)\{1/\Gamma(1+C)-1\}$ is divisible by $C$ as a
holomorphic germ.  Dividing that difference by $C$ and adding
$H/\Gamma(1+C)$ defines $J$.  Uniqueness follows first for $C\ne0$
and then by holomorphy.  The initially convergent Mellin calculation
and uniqueness of meromorphic continuation justify every equality.
\end{proof}

On a ray, the two pole terms become
\[
 -\frac{c}{a+c}\frac{\Gamma(1-at)}{\Gamma(1+ct)}\zeta(bt-1)
 -\frac{c}{b+c}\frac{\Gamma(1-bt)}{\Gamma(1+ct)}\zeta(at-1).
\]
These are holomorphic at zero for $(a,b,c)\in\mathscr D$,
including $c=0$, which proves the removability assertion in
Theorem~\ref{tc:main}.

\subsection{The two determined second derivatives of the remainder}

All derivatives of $J$ in this subsection are at $(0,0,0)$.
Let
\[
 U=J_{AA}=J_{BB},\quad V=J_{AB},\quad
 X=J_{AC}=J_{BC},\quad Y=J_{CC},\quad M=U+2X.
\]
Two elementary restrictions determine $M$ and $Y$ exactly.
Counting positive pairs with a prescribed sum gives
\[
 T(0,0,C)=\zeta(C-1)-\zeta(C).
\]
Consequently \eqref{tc:decomposition} yields
\[
 C J(0,0,C)=\zeta(C-1)-\zeta(C)
                   -\frac{1}{6\Gamma(1+C)}-\frac14.
\]
Taking the coefficient of $C^3$ gives
\begin{equation}
 Y=\frac{r-q}{3}
       -\frac{\gamma^3-3\gamma\zeta(2)+2\zeta(3)}{18}.
 \label{tc:Y}
\end{equation}

For the other restriction, rearrangement of an absolutely convergent
double sum, followed by continuation, proves the classical
quasi-shuffle identity
\begin{equation}
 T(A,0,C)+T(C,0,A)=\zeta(A)\zeta(C)-\zeta(A+C).
 \label{tc:stuffle}
\end{equation}
The polar terms in this sum have the holomorphic continuation
\begin{align}
 \mathscr P(A,0,C)+\mathscr P(C,0,A)
 ={}&\frac{\Gamma(1-A)\Gamma(1-C)}{12}
       \frac{\sin(\pi A)+\sin(\pi C)}{\pi(A+C)}\notag\\
 &-\frac{\zeta(A-1)}{\Gamma(1+C)}
  -\frac{\zeta(C-1)}{\Gamma(1+A)}.
 \label{tc:polar-symmetric}
\end{align}
The quotient in the first line is entire locally, since
$\sin(\pi A)+\sin(\pi C)
=2\sin(\pi(A+C)/2)\cos(\pi(A-C)/2)$.
The coefficient of $A^2C$ in $C J(A,0,C)+A J(C,0,A)$ is $M/2$.
The corresponding coefficient in the first line of
\eqref{tc:polar-symmetric} is
$\gamma(\gamma^2+\zeta(2))/24$.  Coefficient comparison in
\eqref{tc:stuffle} therefore gives
\begin{equation}
 M=p-q+\gamma\zeta''(-1)
      +(\gamma^2-\zeta(2))\zeta'(-1)
      -\frac{\gamma^3+\gamma\zeta(2)}{12}.
 \label{tc:M}
\end{equation}
This calculation uses the symmetric two-variable identity before
specialization; no ambiguous origin value of an individual term is
introduced.

\subsection{Proof of the complete cubic formula}

Let $B(a,b,c)$ be the third derivative at zero of
$\zeta(at)\zeta(bt)+\mathscr P(at,bt,ct)$.  Taylor expansion of
the holomorphic remainder gives exactly
\begin{align}
 W_{a,b,c}'''(0)
 &=B(a,b,c)+3c(a^2+b^2)U+6abcV
                  +6c^2(a+b)X+3c^3Y\notag\\
 &=B(a,b,c)+3\Phi(a,b,c)U+6abcV
                  +3c^2(a+b)M+3c^3Y.
 \label{tc:third-before}
\end{align}
For transparent coefficient checking put $z_j=\zeta^{(j)}(-1)$.
The ordinary Gamma expansion gives
\begin{align*}
 B(a,b,c)={}&-\frac{a^3+b^3}{2}q+3ab(a+b)p
       -c\left(\frac{b^3}{a+c}+\frac{a^3}{b+c}\right)r\\
 &-3c\gamma(a^2+b^2)z_2
       -3c\{\gamma^2(2ab+c(a+b))
                         +\zeta(2)(2ab-c(a+b))\}z_1\\
 &+\frac c{12}\Bigl[
       \gamma^3\{(a+c)^2+(b+c)^2\}
       +3\gamma\zeta(2)(a^2+b^2-2c^2)\\
 &\hspace{38mm}+2\zeta(3)(a^2+b^2-c(a+b)+2c^2)\Bigr].
\end{align*}
For example this follows directly by differentiating
\[
 \log\frac{\Gamma(1-at)}{\Gamma(1+ct)}
  =\gamma(a+c)t+\frac{\zeta(2)}2(a^2-c^2)t^2
                    +\frac{\zeta(3)}3(a^3+c^3)t^3+O(t^4).
\]
Substitute \eqref{tc:Y} and \eqref{tc:M} into
\eqref{tc:third-before}.  Collecting the coefficients of
$\Phi$ and $abc$ gives
\begin{equation}
 W_{a,b,c}'''(0)
 =\Phi(a,b,c)\kappa
   +abc\{6V-6(\gamma^2+\zeta(2))\zeta'(-1)\}
   +\mathscr Q(a,b,c).
 \label{tc:third-two}
\end{equation}
The diagonal has $\Phi(1,1,1)=0$ and
$\mathscr Q(1,1,1)=12p-8q$.  Hence its specialization identifies
the expression in braces as $\Omega-12p+8q$, proving
Theorem~\ref{tc:main}.

An optional independent spectral description of the second constant is
\begin{equation}
 \kappa=-\frac12W_{1,0,2}'''(0)
                +6p-\frac{41}{4}q+\frac72r.
 \label{tc:kappa-Euler}
\end{equation}
Here $T(t,0,2t)=Z(2t,t)$ in the strict Euler double-zeta convention.
Formula~\eqref{tc:kappa-Euler} is a bridge, not a claimed evaluation
of that derivative by ordinary zeta constants.

\subsection{All formal cancellations and an explicit exact identity}

\begin{corollary}[Finite combinations cancelling both cubic coordinates]
\label{tc:cancellation}
Let $\{(a_j,b_j,c_j)\}_{j=1}^N\subset\mathscr D$ and
$w_1,\ldots,w_N\in\mathbb C$.  If
\begin{equation}
 \sum_jw_j a_jb_jc_j=0,\qquad
 \sum_jw_j\Phi(a_j,b_j,c_j)=0,
 \label{tc:criterion}
\end{equation}
then
\begin{equation}
 \sum_jw_jW_{a_j,b_j,c_j}'''(0)
       =\sum_jw_j\mathscr Q(a_j,b_j,c_j).
 \label{tc:cancelled}
\end{equation}
Conditions~\eqref{tc:criterion} are necessary and sufficient for
formal cancellation of the two coordinates in
Theorem~\ref{tc:main}.  Necessity here concerns treating those
coordinates as independent symbols; it is not an arithmetic
independence theorem for their actual values.
\end{corollary}

The locus $\Phi=0$ supplies especially short two-term identities.
For $c\ne0$ it is $a^2+b^2=c(a+b)$.  Its asymmetric positive
integer point $(3,6,5)$ gives
\begin{equation}
 \boxed{\displaystyle
 W_{3,6,5}'''(0)-90W_{1,1,1}'''(0)
   =81p-\frac{403}{2}q-\frac{245}{11}r.}
 \label{tc:two-ray-example}
\end{equation}
The coefficient $90=3\cdot6\cdot5$ cancels the diagonal
coordinate, while $\Phi(3,6,5)=0$ cancels the other one.

For a fixed triple $(a,b,c)$, let
$\phi_c=\Phi(a,b,c)$, $\phi_a=\Phi(b,c,a)$,
$\phi_b=\Phi(c,a,b)$.  A convenient family of weights for the three
cyclic rays is
\[
 (w_c,w_a,w_b)=(\phi_a-\phi_b,\phi_b-\phi_c,\phi_c-\phi_a).
\]
Both conditions in \eqref{tc:criterion} hold identically.
At $(a,b,c)=(1,2,3)$ these weights, divided by four, are $(1,4,-5)$.
Thus we obtain the concrete ordinary-zeta identity
\begin{equation}
 \boxed{\begin{aligned}
 &W_{1,2,3}'''(0)+4W_{2,3,1}'''(0)-5W_{3,1,2}'''(0)\\
 &\qquad=99\zeta'(0)\zeta''(0)
       +\frac{207}{2}\zeta'''(0)
       +\frac{162}{5}\zeta'''(-1).
 \end{aligned}}
 \label{tc:example}
\end{equation}
Every ray in this example has positive coordinates, so no zero-order
factor or restriction to a singular hyperplane is involved.
Equation~\eqref{tc:example} proves a new explicit cancellation relative
to the inspected ProveIt material; it does not evaluate $\Omega$.

\subsection{The exact ambiguity of these local identities at every order}

The appearance of two coordinates at third order has an algebraic
explanation that extends to all orders.  This is a statement about
what the preceding local restrictions determine, not about how many
arithmetically independent constants the actual Tornheim function has.

\begin{theorem}[Holomorphic ambiguity kernel]
\label{tc:kernel}
Suppose two meromorphic germs $T_1,T_2$ have the same polar part
\eqref{tc:polar}, in the precise sense that $T_i-\mathscr P$ is
holomorphic for $i=1,2$. Suppose both germs satisfy symmetry in their first two variables, have
the restriction $T(A,B,0)=\zeta(A)\zeta(B)$ away from the polar set,
and satisfy \eqref{tc:stuffle}.  Their difference has a unique form
\begin{align}
 \Delta T(A,B,C)={}&ABC\,S(A,B,C)\notag\\
 &+AC(A-C)R(A,C)+BC(B-C)R(B,C),
 \label{tc:kernel-form}
\end{align}
where $S$ is a holomorphic germ symmetric in $A,B$ and $R$ is a
holomorphic germ symmetric in its two variables.  Conversely every
expression in \eqref{tc:kernel-form} preserves these local identities.

For homogeneous degree $d\ge3$, the formal ambiguity space has
dimension
\begin{equation}
 \left\lfloor\frac{d^2-1}{4}\right\rfloor.
 \label{tc:kernel-dimension}
\end{equation}
It vanishes in degrees zero, one and two.  Fixing the entire diagonal
germ imposes $S(t,t,t)=0$, and reduces the degree-$d$ dimension by one.
\end{theorem}

\begin{proof}
The common polar term cancels.  The common restriction at $C=0$
then implies $\Delta T=C K(A,B,C)$ for a holomorphic germ $K$
symmetric in $A,B$.  Put $u(A,C)=K(A,0,C)$.  The difference of
\eqref{tc:stuffle} says
\[
 C u(A,C)+A u(C,A)=0.
\]
At $A=0$ this gives $u(0,C)=0$, so $u(A,C)=A v(A,C)$.
After division by $AC$ the identity becomes
$v(A,C)=-v(C,A)$.  An antisymmetric holomorphic germ vanishes on
$A=C$ and is uniquely divisible by $A-C$; its quotient is symmetric.
Thus $u(A,C)=A(A-C)R(A,C)$ for a unique symmetric germ $R$.
The remainder
\[
 K(A,B,C)-u(A,C)-u(B,C)
\]
vanishes at both $A=0$ and $B=0$.  It is therefore uniquely divisible
by $AB$, with a holomorphic quotient symmetric in $A,B$.  This
quotient is $S$ and proves \eqref{tc:kernel-form}.  Direct substitution
proves the converse and also uniqueness.

In homogeneous degree $d$, $S$ and $R$ have degree $d-3$.
Writing $m=d-3$, a symmetric polynomial in two selected variables
and a third unrestricted variable has
$\lfloor(m+2)^2/4\rfloor$ coefficients; a symmetric polynomial in
two variables of total degree $m$ has $\lfloor m/2\rfloor+1$
coefficients.  Their sum is $\lfloor(d^2-1)/4\rfloor$.
On the diagonal the $R$ terms vanish and
$\Delta T(t,t,t)=t^3S(t,t,t)$.  At each homogeneous degree evaluation
of $S$ at $(1,1,1)$ is a nonzero linear functional, so it removes
exactly one dimension.
\end{proof}

At $d=3$ the two free polynomials in \eqref{tc:kernel-form} are
constants; the two direction functions are exactly $abc$ and
$\Phi(a,b,c)$.  This proves that the cancellation criterion is the
complete one obtainable from these local identities alone.  At
degrees $4,5,6$ the corresponding formal dimensions are $3,6,8$.
Additional functional relations may reduce the arithmetic data
further; this theorem does not prejudge that question.

\paragraph{Verification and next questions.}
The accompanying exact symbolic calculation checks the coefficient
collection, the three-ray identity, and the formal kernel dimensions
independently by linear algebra.  A separate numerical script computes
the ray restrictions from the convergent Jonqui\`ere--Mellin formula
and an incomplete-Gamma tail.  Its floating-point residuals are
diagnostics and are not used in the proofs.  At 50 working decimal
digits, the independent comparison for \eqref{tc:example} gives
an absolute residual below $6.3\times10^{-33}$; both sides are
approximately $-451.3535195684862560644674145675405$.
Concrete next questions
are to identify additional functional relations constraining $R$,
to compute the full fourth-order formula in a similarly normalized
basis, and to determine whether the two actual cubic constants admit
arithmetic relations beyond the formal cancellation proved here.
